\documentclass[../main.tex]{subfiles}
\begin{document}

\section{Introduction}

A backdoored classifier answers correctly on ordinary inputs and answers with the attacker's class whenever a trigger is present. Input-level detection asks, for a deployed model and a suspicious input, whether that input carries a trigger, with no retraining and a small set of clean images as the only extra resource. Prediction Shift Backdoor Detection, PSBD \citep{li2025psbd}, answers by switching dropout on at inference and measuring how far the confidence in the model's own prediction falls. A clean input loses a lot, a triggered input barely moves, and the threshold is a quantile of the clean set. The method was built and explained on ResNet-18 and VGG16, with dropout placed after every residual add, and its authors leave the transformer case as future work.

The port to a Vision Transformer is where the question of this paper starts, because a transformer block has more places a perturbation can go than a residual basic block has. A ViT block reads the residual stream twice, once through a LayerNorm into attention and once through a second LayerNorm into the MLP, and writes back twice. Dropout after the residual add perturbs the stream every later block reads, which on \VitBlocks{} blocks compounds far harder than on the few adds of ResNet-18, so the ConvNet rate grid lands past the useful range. The project's founding intuition was that dropout on each branch's output, before the add, would fix that. It does not, and finding out why turned the port into a search over the position of the perturbation among \PositionCount{} named boundaries in the block and the operator applied there among \OperatorCount{} kinds, run over a panel of \PanelCellsTotal{} declared cells of which \PanelCellsClearing{} implanted.

Our contributions follow. First, a deployable configuration. Masking whole tokens at the input of each attention block, at PSBD's own rate rule, beats the published placement by \HeadlineGainAdaptiveAuroc{} mean AUROC on the panel, paired within cell, with the gain positive on every dataset, every rate and every leave-one-attack-out refit, and with a seed spread an order of magnitude below the published placement's. Second, a methodological rule. No comparison across placements is valid at a shared nominal rate, because the same rate is a different disturbance at every position, and every number here is read at a matched clean-validation shift ratio. That rule refuted the founding claim, withdrew an earlier headline of this project, and dissolved the family split that replaced it. Third, a mechanism for the winning placement. A LayerNorm following the injection point absorbs additive noise and not masking, so the operator ranking depends on the position, and attention is the only sublayer that can route a local trigger from its own patches to the class token, which causal patching locates entirely in those patches through the first half of the network. Fourth, a direct test of the published explanation. The prediction shift toward the target class exists on ViT for local triggers on few-class datasets, is absent for global triggers and reads \ShiftToTargetHeadlineCell{} on CIFAR-100 BadNets at \PanelRateLow{}, while detection works in every regime, so the neuron bias account does not carry the method here.

This is a preliminary long version. Every number in it is a macro written by a generator from the results tree, every finding carries a criticality grade in the margin, and the chapters say plainly which results rest on a single checkpoint, which on a single seed, and which are pending. Nothing measured has been left out because it was inconvenient, and the last chapters list what each headline claim still lacks.

\end{document}
